\chapter{Algoritmo de Fruchterman--Reingold para visualización de redes}\label{anexo5}

El algoritmo de Fruchterman--Reingold es un algoritmo heurístico diseñado en 1991 por los científicos de la computación Thomas M. Fruchterman y Edward M. Reingold para dibujar redes mediante fuerzas de atracción y repulsión. El algoritmo simula un sistema físico en el que los vértices representan partículas atómicas que ejercen fuerzas de atracción y repulsión entre sí basadas en sus cargas eléctricas y la fuerza nuclear. Comienza con una distribución aleatoria de los vértices en el plano y después ajusta sus posiciones de manera iterativa, considerando las fuerzas de repulsión y atracción entre ellos. El objetivo es que los nodos conectados por una arista se atraigan entre sí y que, a su vez, todos los vértices se repelan. La proximidad entre los vértices dependerá de su número y el espacio de visualización disponible. Formalmente, el algoritmo busca un estado de equilibrio que minimice la energía total del sistema. Los criterios que guiaron su diseño son distribuir los vértices de forma uniforme en el espacio, minimizar los cruces de aristas, uniformar la longitud de las aristas, reflejar la simetría interna de la gráfica y adaptarse a una delimitación específica \parencite{Fruchterman1991}.  